%------------------------------------------------------------------------------%
\documentclass[fleqn,utf8,aspectratio=169,12pt]{beamer}

\usepackage{gaal}

\title{Equações e Funções}

%------------------------------------------------------------------------------%
\begin{document}

\addtitle

%------------------------------------------------------------------------------%
\section{Coordenadas Cartesianas}

%------------------------------------------------------------------------------%
\begin{frame}{Sistema de Coordenadas Cartesianas}
  \begin{enumerate}[<+->]
  \setlength{\itemsep}{1em}
    \item Cada ponto do plano é associado a um par de números reais $(x, y)$
    \item Cada ``figura'' no plano corresponde a um subconjunto de $\R^2$
    \item Na orientação convencional, $x$ aponta para a direita e $y$ para cima
  \end{enumerate}
\end{frame}

%------------------------------------------------------------------------------%
\begin{frame}{Sistema de Coordenadas Cartesianas}
  \centering
  \includegraphics{B_01_plano_cartesiano}
\end{frame}

%------------------------------------------------------------------------------%
\section{Funções}

%------------------------------------------------------------------------------%
\begin{frame}{Função}
  Função é uma relação que associa a cada elementos
  do domínio um único elemento do contradomínio
\end{frame}

%------------------------------------------------------------------------------%
\begin{frame}{Notação de Função}
  \begin{align*}
    f\colon A & \to B \\[-2mm]
            x & \to f(x)
  \end{align*}
  \pause
  $f$
  \tabto{10mm}
  é o nome da função

  \pause
  $A$
  \tabto{10mm}
  é o conjunto domínio

  \pause
  $B$
  \tabto{10mm}
  é o conjunto contradomínio

  \pause
  $x$
  \tabto{10mm}
  é um elemento qualquer de $A$

  \pause
  $f(x)$
  \tabto{10mm}
  é o único elemento de $B$ associado a $x$
\end{frame}

%------------------------------------------------------------------------------%
\begin{frame}{Gráfico de uma Função Real}
  Conjunto de pontos $(x, y) \in \R^2$ tais que \(y = f(x)\)

  \pause
  \bigskip
  Para cada $x$, no domínio, existe \emph{um único} $y$
\end{frame}

%------------------------------------------------------------------------------%
\framedgraphic{Gráfico de uma Função Real}{B_01_grafico_funcao_real}{}
\framedgraphic{Gráfico de uma Função 2D}{mapa_meteorologico}{}
\framedgraphic{Gráfico de uma Função 3D}{exemplo-peca-tensao}{}

%------------------------------------------------------------------------------%
\section{Equações}

%------------------------------------------------------------------------------%
\begin{frame}{Soluções de uma Equação}

  Qualquer expressão da forma
  \[
    G(x, y) = 0
  \]

  \pause
  Por exemplo,
  \[
    x^2 + y^2 = 4^2
  \]

  \pause
  Queremos \emph{todos os pares} $(x, y)$ que satisfazem a equação

  \pause
  \medskip
  Se faltar, ou sobrar, um, a solução está incompleta e portanto errada
\end{frame}

%------------------------------------------------------------------------------%
\begin{frame}{Soluções de uma Equação}
  \centering
  \includegraphics{B_01_grafico_equacao}
\end{frame}

%------------------------------------------------------------------------------%
\section{Inequações}

%------------------------------------------------------------------------------%
\begin{frame}{Desigualdades}
  Desigualdades geralmente representam regiões no plano (ou espaço)
\end{frame}

%------------------------------------------------------------------------------%
\begin{frame}{Desigualdades}
  \centering
  \includegraphics{B_01_grafico_desigualdade}
\end{frame}

%------------------------------------------------------------------------------%
\section{Linearidade}

%------------------------------------------------------------------------------%
\begin{frame}{Linearidade}
  A \emph{Álgebra Linear} é a área a Matemática que estuda ``objetos'' lineares

  \pause
  \bigskip
  Dizemos que algo é linear quando:
  \begin{itemize}
    \item podemos somar
    \item podemos multiplicar por um escalar (um número)
  \end{itemize}
\end{frame}

%------------------------------------------------------------------------------%
\begin{frame}{Equação Linear}
  Uma equação da forma
  \[
    L(x) = 0
  \]
  é linear se:
  \bigskip

  \pause
  A soma de duas soluções quaisquer também é solução

  \pause
  \quad Se $\alpha$ e $\beta$ são soluções de $L(x)=0$ então $\alpha+\beta$ também é

  \pause
  Um múltiplo de uma solução de também é solução

  \pause
  \quad Se $\beta$ é solução e $a$ é um número então $a\beta$ também é solução
\end{frame}

%------------------------------------------------------------------------------%
\begin{frame}{Exemplo -- Equação Linear}
  A equação
  \[
    L(x, y) = 4x + 2y = 0
  \]
  \emph{é linear}

  \pause
  Note que
  \[
    \alpha = (1, -2)
    \quad\text{e}\quad
    \beta = (-2, 4)
  \]
  são soluções,
  \pause
  então
  \[
    \alpha + \beta
    \pause
    = (1, -2) + (-2, 4)
    \pause
    = (-1, 2)
  \]
  \pause
  também é solução
  \[
    \pause
    L(-1, 2)
    \pause
    = 4(-1) + 2\times 2
    \pause
    = -4 + 4
    \pause
    = 0
  \]
\end{frame}

%------------------------------------------------------------------------------%
\begin{frame}{Exemplo -- Equação Não Linear}
  A equação
  \[
    M(x, y) = 4x^2 + 2y = 0
  \]
  \emph{não é linear}

  \pause
  Note que
  \[
    \alpha = (1, -2)
    \quad\text{e}\quad
    \beta  = (-1, -2)
  \]
  são soluções,
  \pause
  porém
  \[
    \alpha + \beta
    \pause
    = (1, -2) + (-1, -2)
    \pause
    = (0, -4)
  \]
  \pause
  não é solução
  \[
    \pause
    M(0, 4)
    \pause
    = 4\times 0 + 2\times 4
    \pause
    = 0 + 8
    \pause
    = 8
    \pause
    \neq 0
  \]
\end{frame}

%------------------------------------------------------------------------------%
\begin{frame}{Como Resolver}
  Saber se uma equação é linear, ou não, determina como resolvê-la
\end{frame}

%------------------------------------------------------------------------------%
%\section{Lista Mínima}
%
%%------------------------------------------------------------------------------%
%\begin{frame}{Lista Mínima}
%  \begin{enumerate}
%    \item
%  \end{enumerate}
%
%  \vfill
%  Atenção: A prova é baseada no livro, não nas apresentações
%\end{frame}

%------------------------------------------------------------------------------%
\end{document}
%------------------------------------------------------------------------------%
